\documentclass[12pt]{article}
\usepackage{graphicx}
\usepackage{amsmath}
\usepackage{hyperref}
\usepackage{cite}

\title{Advancements in Renewable Energy Technologies: Efficiency, Scalability, and Environmental Impact}
\author{Your Name}
\date{\today}

\begin{document}

\maketitle

\begin{abstract}
The push towards renewable energy has led to significant advancements in solar, wind, and tidal energy technologies. This paper reviews the latest research on these technologies, their efficiency improvements, and scalability challenges. It also critically assesses their environmental impacts and predicts future trends.
\end{abstract}

\section{Introduction}
As concerns over climate change and energy security grow, renewable energy sources have gained prominence. Solar, wind, and tidal energy are pivotal in this transition. This paper discusses current technologies, efficiency improvements, scalability challenges, and environmental impacts, highlighting future trends based on current research.

\section{Solar Energy}
\subsection{Current Technologies}
Recent advancements in solar energy technology include the development of perovskite solar cells and bifacial solar panels\cite{Green2019,Kojima2009}. These technologies exhibit higher efficiencies and lower production costs compared to traditional silicon-based cells.

\subsection{Efficiency Improvements}
Perovskite solar cells have reached efficiencies above 25\% in laboratory settings, rivalling traditional silicon cells\cite{Park2016}. Bifacial panels capture sunlight from both sides, increasing energy yield by up to 30\% in optimal conditions\cite{Sanyo2008}.

\subsection{Scalability Challenges}
Despite these advancements, challenges such as material stability, manufacturing scalability, and recycling need to be addressed to fully realize the potential of these technologies\cite{Yan2019}.

\subsection{Environmental Impact}
Solar energy presents a lower environmental impact compared to fossil fuels. However, issues such as land use and the environmental cost of materials and manufacturing processes need continuous monitoring\cite{Fthenakis2009}.

\section{Wind Energy}
\subsection{Current Technologies}
Modern wind turbines have seen significant advancements in turbine design, materials, and power electronics, which improve their efficiency and reliability\cite{Hau2013}.

\subsection{Efficiency Improvements}
New designs like vertical axis wind turbines (VAWTs) and the integration of AI for predictive maintenance have contributed to efficiency improvements\cite{Islam2013,Brown2019}. The use of larger rotor diameters and higher hub heights have also increased energy capture\cite{Mekonnen2018}.

\subsection{Scalability Challenges}
Challenges include grid integration, variable power output, and the need for extensive transmission infrastructure\cite{Togeby2012}. Offshore wind farms present additional logistical and maintenance difficulties.

\subsection{Environmental Impact}
Wind energy has a relatively low environmental impact. However, concerns remain about noise, visual pollution, and effects on wildlife such as bird and bat populations\cite{Mathews2016}.

\section{Tidal Energy}
\subsection{Current Technologies}
Tidal energy harnesses the energy from ocean currents. Technologies include tidal stream generators and tidal barrages\cite{Khan2009}.

\subsection{Efficiency Improvements}
Advancements in turbine design and materials have increased the efficiency of tidal energy systems. Improvements in computational fluid dynamics (CFD) allow for better design and placement of turbines\cite{Bahaj2011}.

\subsection{Scalability Challenges}
Scalability challenges include high initial costs, maintenance in harsh ocean environments, and potential ecological impacts on marine life\cite{Polagye2010}.

\subsection{Environmental Impact}
Tidal energy systems can impact marine ecosystems. Redirecting water flow can affect tidal ranges and sediment transport, influencing local habitats\cite{Wright2016}.

\section{Conclusion}
Advancements in solar, wind, and tidal energy technologies show promise for improving efficiencies and reducing environmental impacts. However, scalability and ecological considerations present significant challenges. Future research should focus on overcoming these hurdles to maximize the potential of renewable energy sources.

\bibliographystyle{ieeetr}
\bibliography{prompt15_4o}

\end{document}